\documentclass[11pt]{article}
\usepackage[T1]{fontenc}
\usepackage[utf8]{inputenc}
\usepackage{lmodern}
\usepackage[margin=1in,headheight=14pt]{geometry}
\usepackage{amsmath,amssymb,amsthm,mathtools}
\usepackage{microtype,booktabs,tabularx,array,xcolor}
\usepackage{enumitem,listings,fancyhdr,hyperref}
\definecolor{ink}{HTML}{183344}
\definecolor{accent}{HTML}{176B78}
\definecolor{pale}{HTML}{EFF5F7}
\hypersetup{colorlinks=true,linkcolor=accent,citecolor=accent,urlcolor=accent,
 pdftitle={Ramified Semantic Worlds for BusyLean},pdfauthor={Research proposal for Vladimir Reshetnikov}}
\pagestyle{fancy}
\fancyhf{}
\fancyhead[L]{\small\color{ink}BusyLean \textbar\ Ramified semantic worlds}
\fancyhead[R]{\small Research proposal}
\fancyfoot[C]{\thepage}
\setlength{\parskip}{0.35em}
\setlength{\emergencystretch}{1.5em}
\setlength{\parindent}{1.2em}
\setcounter{tocdepth}{2}
\newtheorem{theorem}{Theorem}[section]
\newtheorem{lemma}[theorem]{Lemma}
\newtheorem{corollary}[theorem]{Corollary}
\theoremstyle{definition}
\newtheorem{definition}[theorem]{Definition}
\theoremstyle{remark}
\newtheorem{remark}[theorem]{Remark}
\newcommand{\N}{\mathbb N}
\newcommand{\Tree}{\mathcal T}
\newcommand{\Fam}{\mathcal F}
\newcommand{\World}{\mathcal W}
\newcommand{\Score}{\operatorname{Score}}
\newcommand{\Ext}{\operatorname{Ext}}
\newcommand{\Elev}{\operatorname{Elev}}
\newcommand{\OldStage}{\operatorname{Stage}_0}
\newcommand{\oldidx}{I_0}
\newcommand{\code}[1]{\texttt{#1}}
\newcommand{\funiter}[2]{#1^{[#2]}}
\newcommand{\ctx}{\boldsymbol w}
\lstdefinelanguage{Lean}{keywords={import,open,namespace,end,def,abbrev,structure,inductive,where,match,with,fun,let,by,exact,theorem,termination_by,structural},
 sensitive=true,morecomment=[l]{--},morecomment=[s]{/-}{-/},morestring=[b]"}
\lstset{language=Lean,basicstyle=\ttfamily\footnotesize,keywordstyle=\bfseries\color{ink},
 commentstyle=\itshape\color{black!55},columns=fullflexible,keepspaces=true,
 breaklines=true,frame=single,rulecolor=\color{black!20},backgroundcolor=\color{pale},
 xleftmargin=0.4em,xrightmargin=0.4em,
 literate={→}{{$\to$}}1 {∀}{{$\forall$}}1 {≤}{{$\le$}}1 {⟨}{{$\langle$}}1 {⟩}{{$\rangle$}}1
 {ℕ}{{$\N$}}1 {≠}{{$\ne$}}1 {∧}{{$\wedge$}}1}

\begin{document}
\begin{center}
{\small\sffamily\color{accent}CONSTRUCTIVE LARGE NUMBERS \quad / \quad BUSYLEAN}
\vspace{0.7em}

{\LARGE\bfseries\color{ink}Ramified Semantic Worlds\par}
\vspace{0.3em}
{\Large A proposed Contender 7 for BusyLean\par}
\vspace{0.6em}
{\large Whole-interpreter reflection, self-generated tree indices,\par
and an explicit 42-fold tower lower bound\par}
\vspace{0.8em}
Research proposal prepared for Vladimir Reshetnikov\\
October 12, 2026
\end{center}

\begin{abstract}
The incumbent BusyLean construction reflects a finite-type W-language over
numeric oracle families and iterates this reflection along the trees the
language itself defines. We propose a further step: reflect \emph{entire
semantic worlds}. A world supplies a tree-recursion operator accepting an
arbitrary seed, an arbitrary countably branching well-founded tree, and a
natural argument, together with its own program-generated index. A new
language can access the complete operators and indices of a countable family
of previously constructed worlds. Its bounded program maximum and its own
tree index form a stronger world. We ramify this operation along the preceding
world's own index and repeat the resulting world transformation 42 times.
The resulting closed natural-number definition is computable in principle.
Writing $C_6$ for the incumbent and
$A(x)=\funiter{(u\mapsto 2^u)}{x}(x)$, we prove
\[
                 \boxed{\quad \funiter A{42}(C_6)<C_7.\quad}
\]
We also prove relative language domination, numerical non-representability,
and non-representability of the new tree index. The comparisons do not assume
monotonicity of an interpreter in its oracle. Lean definition and proof source,
an independent small witness evaluator, and reproducible audit scripts are
included. The mathematical argument is complete at the stated level; the
Lean source has not been compiled in this authoring environment.
\end{abstract}

\noindent\colorbox{pale}{\parbox{\dimexpr\linewidth-2\fboxsep\relax}{
\textbf{Validation boundary.} Nine Python test methods pass, including typed
witness checks for all 36 child indices visible to the tower certificate at
input 42. No Lean compiler was available. Accordingly, kernel acceptance,
axiom-closure checks, and the 15-second/60-second contest budgets are
\textbf{not certified}. No such success is asserted by this report.}}

\vspace{0.4em}
\noindent\textbf{Audited upstream snapshot:}\;
\texttt{86b360648163f0249eab79539438b6a83dfe8215}.\\
\textbf{Exact comparison scope:} incumbent Contender 6, proposal P5, and the
faithful Contender 5 port. No exact comparison with each separate P1--P4
proposal is claimed.

\newpage
\begingroup
\small
\setlength{\parskip}{0pt}
\tableofcontents
\endgroup
\newpage

\section{The result and its scope}

The objective is not to append a successor to the old contestant, nor to
assume an uncomputable busy-beaver oracle. It is to describe a stronger
constructive mechanism and give a transparent comparison that remains
meaningful even when neither number can be evaluated.

Let $B$ denote the repository's existing seed family, called
\code{BusyLean.p5Family}. Let $w_0$ consist of its existing forest-recursion
operator and its existing index. We construct an operation $\Elev$ on
semantic worlds and set
\begin{equation}\label{eq:main}
 w_{k+1}=\Elev(w_k),\qquad
 S_k(n)=\Score(w_k,n),\qquad
 C_7=S_{42}(42).
\end{equation}
The base is chosen so that $S_0(n)$ is \emph{definitionally} the incumbent's
function \code{BusyLean.finalFamily n}. In particular, $S_0(42)=C_6$.

The main uniform comparison is stronger than a comparison of two constants:
for every $k\in\N$ and every $n\ge7$,
\begin{equation}\label{eq:growth-overview}
                   A(S_k(n))<S_{k+1}(n),
 \qquad A(x)=\funiter{(u\mapsto 2^u)}{x}(x).
\end{equation}
Monotonicity of the elementary function $A$, \emph{not} of any reflection
operator, then proves
\begin{equation}\label{eq:orbit-overview}
 \funiter A{k}(S_0(n))<S_k(n)\qquad(k\ge1,\ n\ge7).
\end{equation}
The bound in the abstract is the instance $k=n=42$.

There is an additional, more structural reason to regard the proposal as an
advance. The new interpreter exposes every lower world's full operator
\[
 (\text{seed},\text{tree},\text{argument})\longmapsto\text{natural number},
\]
not just its value at 42 or its distinguished unary scoring function.
Programs can choose a lower world at a computed index, construct a new tree,
construct a new seed family, and apply that world's recursion operator to
these new arguments. Section~\ref{sec:separation} proves that the new score
eventually exceeds every fixed program in this richer lower-world language.
The explicit tower is a convenient certificate, not a complete description
of the language's expressive reach.

\paragraph{What is and is not claimed.}
The repository already supplies exact comparisons $C_5<C_6$ and $P5<C_6$
\cite{comparison}. These combine with our comparison to yield $C_5<C_7$ and
$P5<C_7$. Its review ranks P1--P4 using informal growth calibrations
\cite{review}; that is not an exact proof of their values at the contest
input. We do not silently promote that ranking into a new theorem.
Similarly, we assign no unproved proof-theoretic ordinal to $C_7$.

The method uses familiar ingredients---typed interpretation, finite maxima,
well-founded recursion, and diagonalization. The proposed contribution is
their specific integration into a world-valued ramification construction
with explicit, checkable comparison certificates. No claim of global
literature novelty is needed for the results below.

\section{What the incumbent actually supplies}\label{sec:incumbent}

The audited repository defines a type of well-founded trees
\[
 \Tree=\operatorname{WTree}(\N),\qquad
 \mathsf{leaf}:\Tree,\qquad
 \mathsf{node}:(\N\to\Tree)\to\Tree.
\]
These are inductive, potentially infinitely branching trees, not arbitrary
possibly ill-founded graphs. A finite program can denote such a tree by
supplying a terminating function for its branches.

The object language has types generated by
\[
 \tau ::= \mathsf{nat}\mid\tau\to\tau\mid W\tau,
\]
with natural-number recursion and W-recursion into every type. The repository
interprets raw terms totally: type mismatches receive specified default
values. Its finite list \code{BigNumber.allTerms n} is the syntax enumeration
used by the incumbent \cite{definitions}.

Put $\Fam=\N\to\N\to\N$. The incumbent reflector $R_0$ satisfies
\[
 R_0(f,n)=n+1+
 \max\left\{\max_{i\le n} f(i,n),\,
             \max_{t\in L_n}\operatorname{evalWith}(f,t)\right\}.
\]
The forest recursion derived from it is
\begin{align*}
 \OldStage(b,\mathsf{leaf},n)&=R_0(b,n),\\
 \OldStage(b,\mathsf{node}(a),n)&=
 R_0\bigl(i\mapsto[m\mapsto\OldStage(b,a(i),m)],n\bigr).
\end{align*}
The existing index $I_0$ joins every tree denoted by the old language,
with an outer branch for enumeration depth and an inner branch for list
position \cite{contender6,reflection}.

Consequently, the incumbent is exactly
\[
                       C_6=\OldStage(B,I_0,42).
\]
Our base world reuses $\OldStage$, $I_0$, and $B$. The new number's dependency
cone is designed not to contain \code{BusyLean.contender\_6},
\code{BusyLean.finalFamily}, or either proposal-level named contender. This
is more than renaming the previous constant: the reused operator accepts
arbitrary seeds and arbitrary trees, and the new interpreter makes those
parameters available to generated programs.

\subsection{Why oracle monotonicity is the wrong comparison principle}

A program may inspect an oracle result and branch on it. Even an elementary
expression such as ``return 100 if the oracle returns zero, otherwise return
zero'' is not monotone in its oracle. The repository itself explicitly warns
against assuming monotonicity of its reflector \cite{reflectionfacts}.

We therefore do not argue that replacing each child score by a larger score
necessarily makes its parent's reflection larger. Our proofs instead use a
single well-typed witness program whose \emph{exact} denotation is known.
Its value is one of the finitely many values over which the new reflector
maximizes. This approach is unaffected by the behavior of the other
programs in the maximum.

\section{Worlds and a finite semantic signature}\label{sec:worlds}

\begin{definition}[Semantic world]
A world is a pair $w=(\operatorname{run}_w,I_w)$, where
\[
 \operatorname{run}_w:\Fam\to\Tree\to\N\to\N,
 \qquad I_w:\Tree.
\]
Its distinguished score is
\[
                  \Score(w,n)=\operatorname{run}_w(B,I_w,n).
\]
The base world is $w_0=(\OldStage,I_0)$.
\end{definition}

The record deliberately has no field containing its distinguished natural
number. It stores the reusable semantic machinery. The type of a runner is
itself a finite object-language type, since $\Fam$, $\Tree$, and $\N$ all
have codes in the existing type grammar.

Fix a family $\ctx:\N\to\World$ of already available worlds, a live numeric
forest oracle $f:\Fam$, and an input $n$. The extension has the following
nine atoms. All types in this table are finite types of the old grammar.

\begin{center}
\small
\begin{tabularx}{\linewidth}{@{}l l X@{}}
\toprule
Atom & Type & Denotation\\
\midrule
\code{zero} & $\N$ & $0$\\
\code{succ} & $\N\to\N$ & Natural successor\\
\code{input} & $\N$ & Current input $n$\\
\code{oracles} & $\Fam$ & Current numeric forest $f$\\
\code{lowerRun} & $\N\to\Fam\to\Tree\to\N\to\N$ & $i\mapsto\operatorname{run}_{\ctx(i)}$\\
\code{lowerIndex} & $\N\to\Tree$ & $i\mapsto I_{\ctx(i)}$\\
\code{seed} & $\Fam$ & Existing seed $B$\\
\code{iterate} & $(\N\to\N)\to\N\to\N\to\N$ & Finite iteration\\
\code{exp} & $\N\to\N$ & $x\mapsto2^x$\\
\bottomrule
\end{tabularx}
\end{center}

The atoms \code{iterate} and \code{exp} are ordinary System T operations.
They are included explicitly to keep the comparison certificate short and
transparent. They do not smuggle in any large contestant value. Removing
these convenience atoms and replacing them by old-language terms would
change syntax-depth constants, not the basic construction.

\subsection{A combinatory extension rather than a second binder mechanism}

The new raw syntax is
\[
 p::=\mathsf{old}(t)\mid\mathsf{atom}(a)\mid\mathsf{app}(p,p).
\]
An old leaf is interpreted as \code{BigNumber.interpret [] t}. Thus the entire
old closed language is retained unchanged. Application uses a dynamic pair
of a type code and a value. If the function has arrow type $\sigma\to\tau$,
its argument is checked against $\sigma$. Equal type codes transport the
argument unchanged; unequal codes supply the specified default value of
$\sigma$. Applying a non-function yields the natural number zero.

We use exact type equality for this \emph{new} application layer. In
particular, a correctly typed tree argument is not reconstructed by the old
language's structural tree cast. This makes the seven-level witness below
literally denote the old runner on the supplied tree. The semantics inside
an old leaf remains the original semantics. This distinction avoids claiming
that a new casting convention is definitionally identical to every old cast.

No new lambda syntax is needed. Closed old lambda terms provide typed
combinators; application of such combinators to semantic atoms expresses
abstraction over the new constants. For example, the standard typed
$K$- and $S$-combinators are available at all finite types. The implementation
therefore avoids any second de Bruijn-level convention or capture-avoidance
translation.

\subsection{The enumeration really is finite}

Let $C_0=[]$, and write $L_r$ for the old finite term list at depth $r$.
With $\mathsf{Atoms}$ the fixed nine-element list, define
\begin{equation}\label{eq:enumeration}
\begin{split}
 C_{r+1}={}&C_r\mathbin{+\!+}\mathsf{atom}(\mathsf{Atoms})
 \mathbin{+\!+}\mathsf{old}(L_{r+1})\\
 &\mathbin{+\!+}
 [\mathsf{app}(p,q):p\in C_r,\ q\in C_r].
\end{split}
\end{equation}
Duplicates are permitted. Every displayed operation constructs a finite
list; there is no search over semantic equality, halting, or arbitrary
functions.

\begin{lemma}[Persistence and application closure]\label{lem:enumeration}
If $p\in C_r$ and $r\le s$, then $p\in C_s$. Every atom belongs to $C_{r+1}$.
If $p,q\in C_r$, then $\mathsf{app}(p,q)\in C_{r+1}$.
\end{lemma}
\begin{proof}
Persistence for one step is the first summand of
\eqref{eq:enumeration}; iterate it. The other assertions are the second and
fourth summands.
\end{proof}

Arbitrary old terms are \emph{not} unit-cost literals. An old leaf must first
occur in $L_{r+1}$. For completeness, assign type costs
$d(\mathsf{nat})=1$, $d(\sigma\to\tau)=1+\max(d\sigma,d\tau)$, and
$d(W\sigma)=1+d\sigma$. Give old variables $\mathsf{var}(j)$ cost $j+2$,
zero and successor cost 1, each one-type recursor or tree constructor cost
one plus its type cost, application cost one plus the maximum subterm cost,
and lambdas and W-recursors cost one plus the maximum of the costs of their
annotations and subterms. Induction on types and terms, together with
monotonicity of the old enumerations, shows that every old term occurs at its
cost and at all larger depths. A further induction on $p$ shows that every
finite new program belongs to some $C_r$. The supplied Lean separation
theorems take the explicit membership certificate $p\in C_r$, so they do not
need this additional completeness lemma in their dependency cone.

\section{Reflecting whole lower worlds}\label{sec:reflection}

Write $E_{\ctx,f,n}(p)$ for the new natural-number evaluator, returning zero
when the resulting value has a non-natural type. Define
\begin{equation}\label{eq:R}
 R_{\ctx}(f,n)=n+1+
 \max\left\{\max_{i\le n}f(i,n),\,
                    \max_{p\in C_n} E_{\ctx,f,n}(p)\right\},
\end{equation}
with the empty program maximum taken to be zero.

\begin{lemma}[Reflection inequalities]\label{lem:reflection}
For all $\ctx,f,n$, one has $n<R_{\ctx}(f,n)$. If $i\le n$, then
$f(i,n)<R_{\ctx}(f,n)$. If $p\in C_n$, then
$E_{\ctx,f,n}(p)<R_{\ctx}(f,n)$.
\end{lemma}
\begin{proof}
Every relevant value is bounded by one of the displayed maxima. Adding
$n+1\ge1$ makes that bound strict. The input bound follows immediately from
the nonnegativity of the maximum.
\end{proof}

For an arbitrary reflector $R:\Fam\to\N\to\N$, define
\begin{align}
 \operatorname{Fold}_R(b,\mathsf{leaf},n)&=R(b,n),\label{eq:fold}\\
 \operatorname{Fold}_R(b,\mathsf{node}(a),n)&=
 R\bigl(i\mapsto[m\mapsto\operatorname{Fold}_R(b,a(i),m)],n\bigr).
 \nonumber
\end{align}
The recursion is structural on the tree. A higher-order program inside $R$
may call the supplied family at dynamically computed indices and arguments,
but every such family member refers to a proper subtree of the current tree.

The distinction between two interfaces is essential. The atom
\code{oracles} accesses the current numeric forest in \eqref{eq:fold}. The
atoms \code{lowerRun} and \code{lowerIndex} access previously constructed
semantic worlds. The latter are available with arbitrary newly constructed
seeds and trees. They are not confined to the child scores currently used by
the numeric forest.

\section{The new language generates its own tree index}\label{sec:index}

For a fixed family $\ctx$, interpret a program with the live numeric oracle
identically zero and the input zero, and extract a value of type $\Tree$,
using the leaf as a default. Call this tree $T_{\ctx}(p)$. Let $C_k[i]$ be the
$i$th list element, with the zero atom as default. Set
\begin{equation}\label{eq:index}
 I_{\ctx}=\mathsf{node}\left(k\mapsto
        \mathsf{node}\left(i\mapsto T_{\ctx}(C_k[i])\right)\right).
\end{equation}
The new world is
\begin{equation}\label{eq:extend}
 \Ext(\ctx)=
 \bigl(\operatorname{Fold}_{R_{\ctx}},\ I_{\ctx}\bigr).
\end{equation}

This is not a circular definition. The evaluator in $T_{\ctx}$ can inspect
only the supplied lower worlds and the fixed atoms. It does not contain
$\Ext(\ctx)$, $R_{\ctx}$, or $I_{\ctx}$ as one of its own semantic atoms.
Only after the lower-world evaluation and reflection mechanisms have been
defined are the new runner and index packaged together.

\begin{theorem}[Strict domination of every generated tree]\label{thm:tree}
Fix $\ctx$, a seed $b$, and $p\in C_k$. There is $N$ such that for all $n\ge N$,
\[
 \OldStage(b,T_{\ctx}(p),n)<\OldStage(b,I_{\ctx},n).
\]
\end{theorem}
\begin{proof}
Because $p$ is a member of the finite list $C_k$, choose a list position $i$
with $C_k[i]=p$. Take $N=k+i$. The old forest-reflection inequality says that
at an argument $n$, every child with index at most $n$ has score strictly
smaller than its parent. Here $i\le n$ and $k\le n$. Apply that inequality
first to the inner node in \eqref{eq:index}, then to the outer node, and
compose the two strict inequalities. No comparison between different
oracles is used.
\end{proof}

\begin{corollary}[Non-representability of the index]\label{cor:tree-diagonal}
For every $p\in C_k$,
\[
                         T_{\ctx}(p)\ne I_{\ctx}.
\]
\end{corollary}
\begin{proof}
Equality would turn Theorem~\ref{thm:tree}, at $n=N$, into a strict
inequality between a natural number and itself.
\end{proof}

This is a constructive tree diagonalization. It does not need an ordinal
notation system or a rank assignment. In an ordinary set-theoretic model one
may also explain it by ranks: the new index has a two-edge path to every
program-generated tree, so its rank is at least the rank of each such tree
plus two. That observation is optional and supplies no numerical
fast-growing-hierarchy calibration.

The new index also contains each lower world's distinguished tree. The
program
\[
             \mathsf{app}(\mathsf{lowerIndex},\overline i)
\]
denotes $I_{\ctx(i)}$ and occurs by depth $i+2$. Thus index reflection is
strictly extending a semantic collection of trees, not merely choosing a
larger natural parameter for a fixed notation.

\section{A seven-level tower certificate}\label{sec:tower}

Let $P(x)=2^x$, and define finite iteration by
\[
 \funiter f0(x)=x,\qquad
 \funiter f{k+1}(x)=\funiter f{k}(f(x)).
\]
Define
\begin{equation}\label{eq:A}
                         A(x)=\funiter P{x}(x).
\end{equation}
For orientation only, $A(0)=0$, $A(1)=2$, $A(2)=16$, and $A(3)=2^{256}$.
For $x\ge1$, $A(x)$ is at least a tower of $x$ twos with top value 1, because
iteration of $P$ is increasing in its starting argument.

\subsection{The witness and its exact type}

Write $\overline i$ for the unary successor program for $i$, and abbreviate
application by juxtaposition. Define
\begin{align*}
 G_i&=\mathsf{lowerRun}\ \overline i\ \mathsf{seed}\
                          (\mathsf{lowerIndex}\ \overline i),\\
 B_i&=G_i\ \mathsf{input},\\
 Q_i&=\mathsf{iterate}\ \mathsf{exp}\ B_i\ B_i.
\end{align*}
The types are respectively $\N\to\N$, $\N$, and $\N$. At any numeric forest
$f$ and input $n$, exact typing gives
\begin{equation}\label{eq:witness}
 G_i(x)=\Score(\ctx(i),x),\qquad
 E_{\ctx,f,n}(Q_i)=A(\Score(\ctx(i),n)).
\end{equation}
In particular, the value is independent of $f$.

The occurrence of $B_i$ twice is intentional. Its value supplies both the
number of exponential iterations and their starting value. It is not merely
a single exponentiation of the lower score.

\begin{lemma}[Membership certificate]\label{lem:witness-depth}
For each $i\in\N$, $Q_i\in C_{i+7}$.
\end{lemma}
\begin{proof}
A direct application-depth calculation, using
Lemma~\ref{lem:enumeration}, gives the following certificates.
\begin{center}
\begin{tabular}{@{}ll@{}}
\toprule
Expression & Certified enumeration depth\\
\midrule
$\overline i$ & $i+1$\\
$\mathsf{lowerRun}\ \overline i$ & $i+2$\\
$\mathsf{lowerRun}\ \overline i\ \mathsf{seed}$ & $i+3$\\
$\mathsf{lowerIndex}\ \overline i$ & $i+2$\\
$G_i$ & $i+4$\\
$B_i$ & $i+5$\\
$\mathsf{iterate}\ \mathsf{exp}$ & $2$\\
$\mathsf{iterate}\ \mathsf{exp}\ B_i$ & $i+6$\\
$Q_i$ & $i+7$\\
\bottomrule
\end{tabular}
\end{center}
Whenever two argument depths differ, persistence raises the smaller one to
the larger before the application-closure rule is used. No evaluation of
$C_{i+7}$ is involved.
\end{proof}

\begin{theorem}[Whole-world tower bound]\label{thm:extend-tower}
For every $\ctx$, every seed $b$, every tree $a$, and all $i,n$ satisfying
$i+7\le n$,
\[
 A(\Score(\ctx(i),n))<
          \operatorname{Fold}_{R_{\ctx}}(b,a,n).
\]
In particular,
\[
 A(\Score(\ctx(i),n))<\Score(\Ext(\ctx),n).
\]
\end{theorem}
\begin{proof}
By the membership certificate and persistence, $Q_i\in C_n$. Reflection
strictly exceeds its value. Equation~\eqref{eq:witness} identifies that value
for \emph{every} numeric forest. At a leaf the relevant forest is $b$; at a
node it is the family of recursively computed child stages. These two cases
are precisely \eqref{eq:fold}. Finally take $b=B$ and $a=I_{\ctx}$.
\end{proof}

The arbitrary-tree and arbitrary-seed quantifiers matter. They allow the
next construction to choose a complicated self-generated index without
introducing any monotonicity obligation.

\section{World-valued ramification and the final number}\label{sec:ramify}

\begin{definition}[Anchored ramification]
For a base world $w$, define $H_w:\Tree\to\World$ by structural recursion:
\begin{align*}
 H_w(\mathsf{leaf})&=\Ext(i\mapsto w),\\
 H_w(\mathsf{node}(a))&=\Ext(c),
 \quad c(0)=w,\quad c(j+1)=H_w(a(j)).
\end{align*}
Then set $\Elev(w)=H_w(I_w)$.
\end{definition}

At each ramification node, the new interpreter sees the whole semantic worlds
produced by all its child subtrees. Branch zero is an explicit anchor
retaining the base world. The anchor serves two purposes: it gives a small
uniform comparison certificate, and it prevents the proof from needing to
locate the base inside an enormous program enumeration. It contains an
operator and index, not a named predecessor number.

The following is the central portion of the supplied Lean definition.
\begin{lstlisting}
def ramify (base : World) : Tree → World
  | .leaf => extend (fun _ => base)
  | .node f => extend fun i =>
      match i with
      | 0 => base
      | j + 1 => ramify base (f j)
termination_by structural a => a

def elevate (w : World) : World := ramify w w.index

def levels : Nat → World
  | 0 => baseWorld
  | k + 1 => elevate (levels k)

def contender_7 : Nat := score (levels 42) 42
\end{lstlisting}

\begin{theorem}[Uniform ramification margin]\label{thm:ramify-margin}
For every world $w$, every tree $a$, and every $n\ge7$,
\[
                    A(\Score(w,n))<\Score(H_w(a),n).
\]
\end{theorem}
\begin{proof}
Inspect the root constructor of $a$. If it is a leaf, the lower-world family
is constantly $w$. If it is a node, its anchored family satisfies $c(0)=w$.
In either case apply Theorem~\ref{thm:extend-tower} with $i=0$. The theorem
requires no inductive comparison of the recursively generated worlds.
\end{proof}

Substituting $a=I_w$ gives the same inequality for $\Elev(w)$, and hence
\eqref{eq:growth-overview} for every stratum.

\subsection{The elementary monotonicity that is actually used}

\begin{lemma}\label{lem:Amono}
The function $P(x)=2^x$ is nondecreasing and satisfies $x\le P(x)$.
Consequently $A$ is nondecreasing and $x\le A(x)$.
\end{lemma}
\begin{proof}
For the inflationary inequality, induction gives $1\le2^x$ and
\[
             x+1\le2^x+2^x=2^{x+1}.
\]
The same doubling recurrence gives monotonicity of $P$. Induction on $r$
shows that $\funiter P r$ is nondecreasing in its starting value. Since
$P(x)\ge x$, this iterate is also nondecreasing in $r$.
If $a\le b$, then
\[
 \funiter P a(a)\le\funiter P a(b)\le\funiter P b(b),
\]
which is the required monotonicity of $A$. Inflationarity follows by repeated
application of $P(x)\ge x$, including the zero-iteration case.
\end{proof}

\begin{theorem}[The explicit 42-fold bound]\label{thm:main}
For every $n\ge7$ and $k\ge1$,
\[
                  \funiter A{k}(S_0(n))<S_k(n).
\]
In particular, $\funiter A{42}(C_6)<C_7$.
\end{theorem}
\begin{proof}
Let $b_0=S_0(n)$ and $b_{r+1}=A(b_r)$. First prove $b_r\le S_r(n)$ by
induction. Equality holds for $r=0$. At the successor step,
\[
 b_{r+1}=A(b_r)\le A(S_r(n))<S_{r+1}(n),
\]
using Lemma~\ref{lem:Amono} and Theorem~\ref{thm:ramify-margin}. The displayed
successor argument is already strict, proving the first assertion. Put
$k=n=42$ and use $S_0(42)=C_6$.
\end{proof}

This is not a fixed-height exponentiation bound. Even the first application
of $A$ uses the incumbent's entire value as a tower height, and the theorem
iterates this transformation 42 times. The true construction additionally
reflects much richer lower-world programs at each level and along its
self-generated tree indices.

\subsection{A uniform diagonal family}

Set $D(n)=S_n(n)$. This is a single total computable function, with
$C_7=D(42)$.

\begin{corollary}[Domination of every fixed stratum]
For $n\ge7$ and $k<n$, one has $S_k(n)<D(n)$.
\end{corollary}
\begin{proof}
Each successor step strictly increases the score because
$S_j(n)\le A(S_j(n))<S_{j+1}(n)$. Compose these inequalities from $k$ to $n$.
\end{proof}

The diagonalization is over this explicitly constructed hierarchy of worlds,
not over all total computable functions or all Lean definitions.

\section{Relative language domination and separation}\label{sec:separation}

Call a program \emph{forest-oracle-free} when it contains no occurrence of
the atom \code{oracles}. Such a program may still use \code{lowerRun},
\code{lowerIndex}, \code{seed}, \code{input}, all old-language terms, and
arbitrary applications. It can therefore carry out substantial higher-order
computations involving lower worlds.

\begin{lemma}[Independence of the live forest]
If $p$ is forest-oracle-free, then for all $f,g$,
\[
                    E_{\ctx,f,n}(p)=E_{\ctx,g,n}(p).
\]
\end{lemma}
\begin{proof}
Induct on the application tree of $p$. An old leaf ignores both forests.
Every permitted atom has the same denotation for $f$ and $g$. Application
preserves equality of both evaluated operands.
\end{proof}

\begin{theorem}[Domination of the full lower-world language]\label{thm:lang}
If $p\in C_k$ is forest-oracle-free, then for every $n\ge k$,
\[
               E_{\ctx,0,n}(p)<\Score(\Ext(\ctx),n).
\]
\end{theorem}
\begin{proof}
Persistence puts $p$ in $C_n$. At the root of any evaluation of
$\operatorname{Fold}_{R_{\ctx}}$, reflection strictly exceeds the value of
$p$ in the current numeric forest. The preceding lemma identifies this value
with $E_{\ctx,0,n}(p)$ independently of the root forest. Apply this to seed
$B$ and index $I_{\ctx}$.
\end{proof}

\begin{corollary}[No lower-language representation]\label{cor:lang}
No forest-oracle-free $p\in C_k$ satisfies
\[
       E_{\ctx,0,n}(p)=\Score(\Ext(\ctx),n)\quad\hbox{for all }n.
\]
\end{corollary}
\begin{proof}
Use $n=k$ in Theorem~\ref{thm:lang}. The proposed equality contradicts the
strict inequality.
\end{proof}

There is also a direct local version without forest-oracle-freeness. For a
fixed external forest $f$, no $p\in C_k$ can satisfy
$E_{\ctx,f,n}(p)=R_{\ctx}(f,n)$ for every $n$, by the same contradiction at
$n=k$. The forest-oracle-free theorem is stronger for comparing \emph{world
scores}, because a world's internal root forest differs from an externally
chosen $f$.

These theorems are the substantive distinction between the proposal and a
simple numerical amplification. The next world does not merely evaluate one
chosen formula in lower scores. Its score eventually dominates \emph{every
fixed program} of a language that can call the complete lower semantic
operators on arbitrary program-generated arguments. Theorems are stated with
explicit syntax membership bounds: they do not pretend that a program of
unbounded depth is already dominated at input 42.

\section{Why the number is computable in principle}\label{sec:computability}

A bare statement that a fixed natural number is computable would be nearly
vacuous: each fixed natural has a constant program. What is needed here is an
effective route from the displayed definition to its value, and preferably
a uniform family such as $D$ whose totality is justified.

\begin{theorem}[Constructive totality]
The definitions in Sections~\ref{sec:worlds}--\ref{sec:ramify}, starting from
the repository's constructive base machinery, give a total effective
procedure for $D(n)$ for every $n\in\N$. In particular they define $C_7$ as
a computable natural number.
\end{theorem}
\begin{proof}
The proof follows the construction's dependency order.

First, type-code equality and the new raw interpreter are structurally
recursive procedures on finite codes and finite programs. Old leaves use the
existing total W-language interpreter. Each primitive is a total function
when its supplied lower worlds are effectively available.

Second, $C_n$ is a finite list produced by natural-number recursion. Every
maximum in $R_{\ctx}$ is over that finite list or a finite initial segment of
$\N$. It therefore terminates, however large the list and its values may be.

Third, $I_{\ctx}$ is built using two W-node constructors and total branch
functions. To request a branch, one computes a finite list lookup and the
tree denotation of one terminating program. Constructing the tree does not
require materializing all infinitely many branches.

Fourth, $\operatorname{Fold}_{R_{\ctx}}$ recurses structurally on a W-tree.
Its recursive functional arguments refer only to proper subtrees. Programs
may make enormous numbers of calls and choose their arguments dynamically,
but those calls do not introduce a self-reference to the current tree.

Fifth, $H_w$ recurses structurally on its schedule tree and returns a world.
For a node, every non-anchor lower world comes from a proper child subtree.
The new interpreter sees these lower worlds as already supplied data; it
never interprets its own semantic definition. Finally, $w_n$ is obtained by
natural-number recursion on $n$, and $D(n)$ is its score at $n$.

All operations are therefore combinations of finite recursion, total
interpretation, and structural recursion on inductive W-trees. No halting
oracle, minimization over an undecidable predicate, or classical selection of
a maximizing program occurs.
\end{proof}

\paragraph{Arbitrary parameters versus the closed candidate.}
An arbitrary set-theoretic element of a function type need not come with an
effective presentation. The parametric theorems about $\ctx$ are statements
about total semantic functions. The computability assertion concerns
\emph{effectively supplied} lower worlds, and ultimately the closed family
built from $w_0$. It does not assert that an arbitrary external mathematical
oracle becomes computable merely by appearing in a Lean parameter.

\paragraph{Why there is no diagonal paradox.}
The lower-world language does not contain its new reflector or its new index
as an atom. Corollaries~\ref{cor:tree-diagonal} and \ref{cor:lang} show exactly
why that boundary is essential. A total evaluator can interpret an explicitly
weaker, stratified language without interpreting itself. The finite
hierarchy and its diagonal remain within Lean's host theory; they do not
supply a total interpreter for every total Lean function.

The intended Lean implementations use ordinary \code{def} declarations,
finite lists, inductive types, and structural recursion. This is the standard
termination route described by Lean's language reference \cite{leanrec}.
The theoretical argument is distinct from a claim that a particular source
file has passed the kernel, which requires actually running the toolchain.

\section{Lean artifacts, audit design, and validation}\label{sec:validation}

\subsection{Artifact organization}

The archive contains a standalone Lake project pinned to the audited BusyLean
commit and to Lean 4.34.0, the repository's declared version. The principal
modules are:

\begin{center}
\small
\begin{tabularx}{\linewidth}{@{}l X@{}}
\toprule
File & Role\\
\midrule
\code{BusyLeanNext/Definitions.lean} & Worlds, exact casts, atoms, evaluator,
finite enumeration, reflection, index, ramification, final number, and
comparison-witness syntax.\\
\code{BusyLeanNext/Proofs.lean} & Enumeration membership, exact witness
semantics, tower inequalities, comparisons, language separation, and tree
index non-representability.\\
\code{test/Audit.lean} & Transitive constant-closure and axiom checks.\\
\code{test/Smoke.lean} & Only small primitive equalities and theorem-signature
checks; no contender evaluation.\\
\bottomrule
\end{tabularx}
\end{center}

The proofs specialize general inequalities only at the end. They do not ask
Lean's simplifier to compute a candidate number or expand a bounded maximum
at depth 42. Correctly typed application witnesses are proved at symbolic
worlds and inputs, before any enormous world is substituted.

The supplied audit is intended to reject dependency on the predecessor's
named constant or final-family wrapper, as well as on
\code{Classical.choice}, \code{sorryAx}, or any custom axiom. The allowed
axioms are exactly \code{propext} and \code{Quot.sound}, matching the
repository's rule. That is a specification of the audit, not a statement that
its output has already been obtained.

\subsection{What was actually tested}

Nine Python unit-test methods passed. They verify exact simple-fragment types,
the depth formula $i+7$ for 100 indices, absence of the live numeric oracle in
the witness, correct routing of lower runners and tree identities, exact
small tower values, basic exponential inequalities, exact-cast behavior,
wrong-function defaults, and the arithmetic finite-maximum inequality.
The semantic witness is exercised for every $i\in\{0,\ldots,35\}$ at input
42, with four deliberately small lower-world scores. These 144 semantic
cases ensure that the same indexed world's runner and tree are used and that
the live numeric oracle is not accidentally consulted.

The reference evaluator deliberately does \emph{not} implement or enumerate
the complete old W-language. It is an independent check of the small
certificate and application layer, not a test of $C_7$ and not a substitute
for formal proof. Its safety limits apply only to this testing program;
they are not part of the mathematical definition.

A lexical delimiter check was also performed on the supplied Lean files.
This catches only elementary source damage. It proves nothing about typing,
termination checking, theorem acceptance, or axiom dependencies.

\subsection{What remains unverified}

No Lean executable was present in the authoring environment. Attempts to
obtain the pinned compiler did not provide an executable toolchain. The
included \code{audit/lean-validation.json} therefore records
\code{status: unavailable} and an empty list of executed Lean phases.
There are no fabricated build times or kernel-success logs.

In particular, all of the following are outstanding: acceptance of the
source by Lean 4.34.0; successful kernel checking of the comparison proofs;
the intended transitive dependency and axiom checks; and compliance with the
15-second definition and 60-second proof budgets. The source contains no
placeholder proofs, but it is still \emph{uncompiled formalization source},
not a fully validated contest submission. Elaboration or library-interface
repairs may be necessary. Such repairs must preserve the mathematical
semantics and be followed by the full audit.

To run the reproducibility checks after installing the pinned toolchain:
\begin{lstlisting}[language={}]
python3 scripts/check_lean.py
python3 -m unittest discover -s tests -v
\end{lstlisting}
The Lean script fetches dependencies before timing, removes project and
dependency build outputs while retaining Lean core, builds definitions and
proofs in separate phases, and then runs the audit and smoke tests. It fails
rather than treating a timeout or missing compiler as a pass.

\subsection{The theorem-to-source correspondence}

\begin{center}
\small
\begin{tabularx}{\linewidth}{@{}l X@{}}
\toprule
Mathematical result & Principal Lean declaration\\
\midrule
Witness depth & \code{towerProgram\_mem}\\
Exact tower denotation & \code{evalNat\_towerProgram}\\
Any-seed, any-tree margin & \code{tower\_lt\_foldReflect}\\
Anchored ramification margin & \code{tower\_lt\_score\_ramify}\\
42-fold incumbent bound & \code{towerOrbit\_contender6\_lt\_contender7}\\
Required strict comparison & \code{contender\_6\_lt\_contender\_7}\\
All lower-language programs & \code{lower\_language\_eventually\_lt}\\
No lower-language representation & \code{no\_lower\_language\_representation}\\
No self-index representation & \code{no\_programTree\_eq\_index}\\
Fixed-stratum diagonal domination & \code{fixed\_level\_lt\_diagonal}\\
\bottomrule
\end{tabularx}
\end{center}

\section{Research directions and unresolved comparisons}

\subsection{Complete the kernel and resource audit}

The immediate engineering objective is to compile this exact source at the
pinned version, repair any elaboration mismatches, and publish the complete
logs. Definitions and proofs should remain separated so that proof-library
imports do not inflate the definition budget. Exact casts and explicit
membership certificates should be retained unless their replacements come
with semantic-equivalence proofs.

\subsection{Certify P1--P4 at the fixed input}

A worthwhile next theorem is a source-to-syntax compilation of every earlier
proposal into this language, with a verified enumeration-depth certificate.
That would replace the repository's informal ranking by exact pointwise
comparisons. An asymptotic ordinal estimate is insufficient: eventual
domination of functions does not by itself compare their values at 42.
This is a concrete missing comparison, not a limitation of the incumbent
bound proved here.

\subsection{Remove convenience atoms without losing the 42 budget}

Translate \code{iterate} and \code{exp} into closed old-language terms and
prove semantic equality. Then optimize the witness's exact enumeration
cost. Similar work could remove the explicit zero and successor atoms,
although those are already elementary. The interesting optimization target
is not merely fewer source characters; it is a smaller, independently
verified syntactic certificate with predictable elaboration cost.

\subsection{Replace the inherited seed guard}

The base seed inherits the old P5 guard involving the old-language maximum.
It is constructive and allowed by the present dependency rule, but it remains
a historical comparison device. A fully intrinsic typed embedding of the
original System T contender, with an explicit short program certificate,
could remove this guard while retaining an exact $C_5$ comparison.

\subsection{Analyze world ramification as an operation on function classes}

The language-domination theorem is a precise starting point for a
proof-theoretic analysis. One should identify which classes of functionals
are represented at each stratum, and establish simulations with explicit
size translations before assigning an ordinal calibration. The construction
combines tree recursion and interpreter promotion; it should not be assigned
a growth ordinal by mechanically multiplying the incumbent's informal
estimate.

\subsection{Study observational equivalence and tree transformations}

Exact type casts retain tree values, whereas the old language's self-casts
may reconstruct them extensionally. A logical relation connecting these
representations would clarify compiler-correctness results and permit
aggressive combinatory optimizations. It would also separate genuine
semantic strength from harmless changes in tree presentation or enumeration
order.

\subsection{Higher-order world operators and alternative schedules}

The present number uses 42 repetitions of a world transformation that is
already ramified along its own index. One can investigate languages whose
atoms expose higher-order operations on worlds, or systematically compare
alternative anchored schedules. Such an extension needs a fresh stratification
argument: merely inserting the current interpreter into its own atom table
would violate the separation theorems and destroy the intended totality
boundary.

\subsection{A reusable constructive reflection framework}

Most arguments above depend only on a finite cumulative syntax enumeration,
a total typed evaluator, a W-tree type, and a small set of semantic atoms.
Abstracting these hypotheses could turn the seven-level witness method,
anchored ramification lemma, and self-index separation theorem into reusable
Lean components. Future contestants could then change their object language
without re-proving all low-level finite-maximum facts.

\section{Conclusion}

The proposed number is
\[
                        C_7=\Score(w_{42},42),
 \qquad w_{k+1}=H_{w_k}(I_{w_k}).
\]
Its construction promotes complete lower-level semantic operators into a
new total language, generates that language's tree index, and ramifies the
result at the level of semantic worlds. The explicit certificate proves
\[
                 \funiter A{42}(C_6)<C_7,
 \qquad A(x)=\funiter{(u\mapsto2^u)}{x}(x).
\]
The language and tree separation theorems explain why the mechanism is more
than one preselected numerical amplification. Computability follows from a
stratified dependency graph and structural recursion, not from an oracle or
an unsupported normalization assumption. The remaining validation step is
concrete and falsifiable: compile the supplied Lean source, execute its
dependency audit, and measure it against the repository's resource rules.

\begin{thebibliography}{9}
\raggedright
\bibitem{repository}
V. Reshetnikov and contributors.
\emph{BusyLean}, audited commit
\texttt{86b360648163f0249eab79539438b6a83dfe8215}.
\url{https://github.com/VladimirReshetnikov/BusyLean}.
Source retrieved October 12, 2026 (America/Los\_Angeles).

\bibitem{definitions}
BusyLean contributors.
\emph{P5 W-language definitions},
\path{proposals/P5-w-reflection/lean/BigNumber/Definitions.lean},
at the commit in \cite{repository}.

\bibitem{contender6}
BusyLean contributors.
\emph{Contender 6}, \texttt{BusyLean/Contender6.lean},
at the commit in \cite{repository}.

\bibitem{reflection}
BusyLean contributors.
\emph{Reflection over the W-language},
\texttt{BusyLean/Reflection.lean}, at the commit in \cite{repository}.

\bibitem{reflectionfacts}
BusyLean contributors.
\emph{Inequalities for reflection and forest reflection},
\texttt{BusyLean/ReflectionFacts.lean}, at the commit in \cite{repository}.

\bibitem{comparison}
BusyLean contributors.
\emph{Contender 6 comparison theorems},
\texttt{BusyLean/Comparison.lean}, at the commit in \cite{repository}.

\bibitem{review}
BusyLean contributors.
\emph{Review of the five incoming proposals},
\texttt{proposals/REVIEW.md}, at the commit in \cite{repository}.
Its ordinal-based rankings are explicitly described as informal.

\bibitem{leanrec}
Lean contributors.
\emph{The Lean Language Reference: Recursive Definitions}.
\url{https://lean-lang.org/doc/reference/latest/Definitions/Recursive-Definitions/}.
Consulted October 12, 2026. The executable project is pinned to Lean 4.34.0,
not to a moving ``latest'' toolchain.
\end{thebibliography}
\end{document}
